% ===========================================================================
%  14 — Design system
% ===========================================================================
\chapter{\emph{Design system} e camada de interface}
\label{cap:ds}

\begin{regra}
\textbf{Todo elemento de interface do hub e das experiências vem do
\cd{@unooria/design-system}, sempre na última versão publicada.} Componente de
interface novo se escreve \emph{no repositório do \emph{design system}} e chega
aqui por versão --- não em \arq{packages/ui}, não dentro de uma experiência.
Antes de criar um botão, um \emph{card} ou um campo: procure na seção
correspondente do pacote, ela provavelmente já tem.
\end{regra}

O pacote é npm \textbf{público} no escopo \cd{@unooria} --- sem \emph{registry}
privado, sem \arq{.npmrc} no repositório. A fonte é o repositório
\arq{storybook-unooria} (Storybook em \cd{:6007}, publicado por \emph{tag}
\cd{v*} no GitHub Actions).

\section{Duas famílias no mesmo pacote}

O pacote publica \textbf{seções} independentes, cada uma com \emph{namespace} de
token e prefixo de classe próprios. Elas nunca se cruzam.

\begin{tabularx}{\linewidth}{@{}L{24mm}L{22mm}L{20mm}X@{}}
\toprule
\headrow \thd{Seção} & \thd{Import} & \thd{Tokens} & \thd{Onde vale}\\
\midrule
\cd{unooria} (39) & raiz do pacote &
  \cd{--unooria-*}, classe \cd{.unooria-theme-light|dark} &
  Hub: \emph{landing} (\rota{/}), \emph{login}, \rota{/professor},
  \rota{/admin}, \rota{/analise}; \emph{header} de todos os papéis; telas
  \textbf{pós-revelação} do estudante\\
\zebra \cd{instanooria} (21) & \cd{/instanooria} & \cd{--insta-*}, \cd{insta-} &
  exp-01\\
\cd{unitube} (30) & \cd{/unitube} & \cd{--uni-*}, \cd{uni-} ---
  \textbf{só escuro} & exp-02\\
\zebra \cd{whatsappmobile} (29) & \cd{/whatsappmobile} & \cd{--wa-*} --- só escuro
  & \textbf{exp-07} --- da tela de bloqueio à validação de Amara, inclusive os
  erros diegéticos\\
\cd{chatgpt} (25) & \cd{/chatgpt} & \cd{--gpt-*},
  \cd{data-gpt-theme} & \textbf{exp-03 inteira} --- ela não tem revelação no
  software, então nunca troca para a seção da plataforma\\
\zebra \cd{tiktooria} (35) & \cd{/tiktooria} & \cd{--tik-*} &
  --- (dois \emph{contextos} de superfície: cromo claro de lista e perfil,
  imersivo preto de \emph{feed} e câmera --- não dois temas)\\
\bottomrule
\end{tabularx}

\begin{nota}[title={NOTA --- ``pós-revelação'' virou uma tela de verdade}]
Por muito tempo a linha ``telas pós-revelação do estudante'' descrevia um único
componente: o \cd{DonutChart} do relatório, num \emph{chunk} tardio. O resto do
relatório era Tailwind do clone sobre \cd{bg-night}.

Com a ADR-41 (\ref{adr:relatorio-expandido}) o relatório sai do aparelho e passa
a ser claro e largo, e a autorização deixa de ser teórica. A fronteira não
mudou: quem a marca é a REVELAÇÃO --- sete telas, fundo \cd{night}, sem nada da
plataforma. Antes dela, o disfarce; depois, a língua da escola.

A trava técnica continua a mesma e não pode cair: o CSS do DS entra por
\cd{React.lazy}, no \emph{chunk} do relatório, e nunca no \emph{bundle} do
feed --- lá ele seria legível no DevTools antes da hora.
\end{nota}

\begin{atencao}
A seção \cd{unitube} é \textbf{só escura}: o Figma não publica variante clara, e
\textbf{derivar uma é proibido}. Uma paleta clara inventada não é o clone --- é
uma aproximação, e a experiência mede reconhecimento.
\end{atencao}

\section{A linha que separa as seções é a REVELAÇÃO}

\begin{figure}[H]
\centering
\fitwidth{%
\begin{tikzpicture}[x=1cm, y=1cm]

  \draw[iframe] (0.5,0.7) rectangle (8.0,-4.6);
  \node[iframelbl] at (0.7,0.85) {Dentro do \texttt{iframe} --- antes da revelação};

  \node[c4component, text width=26mm, fill=unooPaper, draw=unooAxis] (feed) at (2.4,-0.6)
    {\emph{feed} de vídeos};
  \node[c4component, text width=26mm, fill=unooPaper, draw=unooAxis] (cromo) at (6.0,-0.6)
    {cromo, busca, menu};
  \node[c4component, text width=26mm, fill=unooPaper, draw=unooAxis] (play) at (2.4,-2.4)
    {\emph{playlists}};
  \node[c4component, text width=26mm, fill=unooPaper, draw=unooAxis] (canal) at (6.0,-2.4)
    {canal do estudante};
  \node[c4component, text width=26mm, fill=unooPaper, draw=unooAxis] (toast) at (4.2,-3.9)
    {\emph{toast}: \enquote{Vídeo salvo!}};
  \node[cxlabel, anchor=north] at (4.2,0.55) {\texttt{@edu/ui/ds-unootube} --- seção do CLONE};

  \draw[c4boundary, draw=unooNavy] (9.0,0.7) rectangle (15.5,-4.6);
  \node[c4boundarylbl, fill=unooSurface, text=unooNavy] at (9.2,0.85)
    {Depois da revelação --- ainda no mesmo \texttt{iframe}};

  \node[c4component, text width=26mm] (rank) at (10.9,-0.6) {ranking};
  \node[c4component, text width=26mm] (d1) at (14.2,-0.6) {diálogo 1 da missão};
  \node[c4component, text width=26mm] (d2) at (10.9,-2.4) {diálogo 2 da missão};
  \node[c4component, text width=26mm] (fim) at (14.2,-2.4) {fechamento};
  \node[cxlabel, anchor=north] at (12.3,0.55) {\texttt{@edu/ui/ds} --- seção da PLATAFORMA};

  \node[actnote, text width=52mm, anchor=north west] at (0.5,-5.4)
    {\textbf{A linha é a revelação, não o arquivo.} Enquanto o estudante acha que
     está num aplicativo de vídeo, tudo é clone. Da virada da missão em diante ele
     está de volta na Unooria, e a linguagem institucional é \emph{parte do que se
     ensina}. Um \texttt{Button} da plataforma no \emph{feed} quebra a ilusão que
     a experiência está medindo.};

  \node[actnote, text width=52mm, anchor=north west] at (9.0,-5.4)
    {\textbf{O \emph{toast} não cruza a linha:} ele fala de dentro da ficção
     (\enquote{Vídeo salvo!}). É o exemplo mais fino da regra --- o critério não é
     \enquote{é depois no tempo?}, é \enquote{quem está falando?}.};

  \draw[c4relhi] (8.0,-2.0) -- (9.0,-2.0) node[cxlabel, midway, above]{revelação};

\end{tikzpicture}}
\caption[A linha da revelação]{A linha da revelação, na exp-02. As duas seções do
  \emph{design system} coexistem no mesmo \emph{bundle} e a fronteira entre elas é
  pedagógica, não técnica.}
\label{fig:linha-revelacao}
\end{figure}

\begin{nota}
\textbf{A linha tem duas pontas, não uma.} A versão 1.1.0 da exp-02 abriu com uma
\textbf{tela de entrada} que explica a missão, e ela usa a seção da
\textbf{plataforma} --- antes do \emph{feed}, ou seja, \emph{antes} da revelação.
Não é exceção à regra: é a regra aplicada pelo critério que ela mesma enuncia,
\enquote{quem está falando?}. No \emph{briefing} quem fala é a escola, e não há
disfarce a proteger --- a criança ainda não viu um \emph{Short}. O disfarce começa
no toque em \enquote{Começar missão}. Ver \adrr{ADR-22}.
\end{nota}

\section{Consumo: sempre pela ponte}

\begin{regra}
\textbf{Consumir sempre pela ponte em \arq{packages/ui}}, nunca importando o
pacote direto num componente:

\begin{tabularx}{\linewidth}{@{}L{28mm}L{22mm}X@{}}
\cd{@edu/ui/ds} & \cd{unooria} & plataforma\\
\cd{@edu/ui/ds-insta} & \cd{instanooria} & exp-01\\
\cd{@edu/ui/ds-unootube} & \cd{unitube} & exp-02\\
\cd{@edu/ui/ds-chatgpt} & \cd{chatgpt} & exp-03\\
\cd{@edu/ui/ds-whatsapp} & \cd{whatsappmobile} & exp-07\\
\end{tabularx}
\end{regra}

A ponte existe por dois motivos, e os dois estão registrados nos arquivos:

\begin{enumerate}
  \item \textbf{Marcar \cd{'use client'} de uma vez}, para o App Router do
    Next.js.
  \item \textbf{Desfazer a colisão de nomes.} \cd{Icon}, \cd{Avatar},
    \cd{Badge}, \cd{Button}, \cd{Logo} e \cd{TextField} existem nas
    \emph{duas} famílias. Os clones só saem pelo \emph{namespace}
    (\cd{InstanooriaUI.Icon}), e \textbf{destruturar na ponte} é o que mantém o
    \emph{tree-shaking} e evita o erro \cd{TS2742} causado pelo
    \cd{@types/react} interno do pacote.
\end{enumerate}

\begin{nota}
Clone novo em uso = arquivo novo \arq{packages/ui/src/ds-<clone>.ts} +
\cd{exports} no \arq{packages/ui/package.json}. Mesmo padrão, sempre.
\end{nota}

\section{CSS: uma vez, no ponto de entrada}

\begin{tabularx}{\linewidth}{@{}L{34mm}X@{}}
\toprule
\headrow \thd{Aplicação} & \thd{Onde o CSS entra}\\
\midrule
Hub & \arq{components/casca-plataforma.tsx} --- e é \emph{por isso} que
  \rota{/exp/*} e \rota{/dev-embed}, que estão fora do grupo de rotas
  \cd{(plataforma)}, também têm o \emph{design system}.\\
\zebra Experiência & O próprio \arq{src/main.tsx}, com o CSS da \emph{sua}
  seção.\\
\bottomrule
\end{tabularx}

\begin{regra}
\textbf{Nunca importe o consolidado \arq{/styles.css} numa experiência}: ele
arrasta as seis seções e a tipografia da plataforma para dentro do clone.
\end{regra}

\subsection{No hub, o CSS não basta: as fontes também}

\begin{atencao}
O \arq{unooria.css} declara \cd{Open Sans} e \cd{Material Symbols Rounded} por
\cd{@import} de URL remota --- e o \textbf{Next remove \cd{@import} externo} de
CSS importado por JS, \emph{silenciosamente}: a regra continua lá, o
\cd{getComputedStyle} responde \enquote{Open Sans} e o navegador desenha com a
fonte do sistema.

Como \cd{Material Symbols Rounded} é fonte de \textbf{ligadura}, o sintoma é o
\textbf{ícone virando texto}: onde deveria haver um ícone de adicionar, aparece
\enquote{addInvestigação}.

Por isso \arq{components/casca-plataforma.tsx} também monta
\cd{FontesDaPlataforma} (\cd{@edu/ui/fontes}). Numa experiência Vite o
\cd{@import} sobrevive e isso não é necessário.
\end{atencao}

\begin{atencao}
\textbf{Segundo caminho para o mesmo sintoma: \cd{icon} como \emph{string}.}
Componentes de navegação do pacote (\cd{SideBar}, \cd{Tab}) tipam \cd{icon} como
\cd{React.ReactNode}, não como nome de ícone. Passando a \emph{string}, o item
renderiza o \textbf{texto cru}: a barra lateral da área administrativa escrevia
\enquote{group}, \enquote{school} e \enquote{upload\_file} na lateral --- com a
mesma cara de fonte que não carregou, e com a fonte perfeitamente carregada.

A forma certa é \cd{icon=\{<Icon name="group" />\}}. Quem pegou foi o
\emph{smoke} visual: \emph{typecheck} não reclama, porque \cd{string} é um
\cd{ReactNode} válido.
\end{atencao}

\begin{regra}
\textbf{Quem importa, declara.} Todo aplicativo que escreve
\cd{@unooria/design-system} no código tem a dependência no seu
\arq{package.json} --- hub incluído, porque ele importa o \arq{unooria.css}.
Depender do \emph{hoisting} do pnpm via \arq{packages/ui} funcionava
\emph{por acaso}: bastava a ponte trocar de versão (ou o \arq{node\_modules} ser
reinstalado com outro \emph{layout}) para o CSS do hub sumir sem ninguém ter
mexido nele.
\end{regra}

\section{Frames: dois níveis, não empilhar}

\begin{tabularx}{\linewidth}{@{}L{34mm}X@{}}
\toprule
\headrow \thd{Componente} & \thd{Onde vale}\\
\midrule
\cd{PhoneFrame}, \cd{DesktopFrame} & Cascas de aparelho \textbf{do clone},
  publicadas pelo pacote. Servem ao Storybook e a telas soltas.\\
\zebra \cd{DispositivoFrame} (hub) & \textbf{Dentro da plataforma, quem desenha o
  aparelho é o hub}, dirigido pelo \cd{viewport} do manifesto. A experiência
  embedada preenche o \cd{iframe} e não traz casca própria.\\
\bottomrule
\end{tabularx}

\section{Tema}

\begin{tabularx}{\linewidth}{@{}L{30mm}X@{}}
\toprule
\headrow \thd{Seção} & \thd{Como recebe o tema}\\
\midrule
Clone (\cd{unitube}) & Lê \cd{data-uni-theme} na raiz do documento.\\
\zebra Plataforma (\cd{unooria}) & Recebe \cd{theme} por \textbf{prop}.\\
\bottomrule
\end{tabularx}

\begin{atencao}
Quem troca o tema escreve nos \textbf{dois} (\arq{exp-02/src/lib/tema.ts}),
senão o estudante escolhe escuro no \emph{feed} e o ranking volta branco --- porque
o ranking é da seção da plataforma, que não lê o atributo da raiz. Na exp-02 a
escolha é do estudante (Perfil \seta{} APARÊNCIA) e é lembrada em
\cd{localStorage}, que a regra 4 proíbe para \emph{token}, não para preferência
de tela.
\end{atencao}

\section{Versão: atualizar é ato deliberado}

\begin{atencao}
\cd{\^{}0.1.11} \textbf{não} pega \cd{0.2.0}: \emph{caret} em \cd{0.x} só sobe
\emph{patch}. \enquote{Última versão} é, portanto, ativo:

\begin{lstlisting}[language=bashunoo]
npm view @unooria/design-system version     # qual e a ultima publicada
pnpm up -r @unooria/design-system --latest  # sobe em ui/hub/exps de uma vez
# commitando o lockfile E os package.json juntos
\end{lstlisting}
\end{atencao}

\section{Compensações locais}

\begin{regra}
\textbf{Componente que falta no pacote não nasce na experiência.} Vira
\emph{compensação local} marcada no CSS (comentário \cd{(b) COMPENSAÇÃO}) e um
pedido ao \arq{storybook-unooria}. A compensação é \textbf{exceção controlada e
decrescente}: cada bloco compensa uma lacuna real, anotada no comentário, e sai
quando o pacote publicar a correção.
\end{regra}

\subsection{Lacunas abertas, verificadas na 0.1.14}

\begin{nota}
As quatro primeiras da tabela abaixo foram \textbf{atendidas na 0.1.17} e as
compensações correspondentes foram apagadas da exp-02. Duas lacunas novas
apareceram na 1.1.0 da experiência e estão registradas logo depois.
\end{nota}

\begin{tabularx}{\linewidth}{@{}L{40mm}X@{}}
\toprule
\headrow \thd{Lacuna} & \thd{Consequência}\\
\midrule
\cd{unitube} não publica \emph{toast}, diálogo, tabela nem estado vazio &
  Quatro compensações locais na exp-02.\\
\zebra \cd{ChannelHeader} não tem \cd{showSubscribe} & O botão de inscrever-se é
  \emph{sempre} desenhado, e o canal da exp-02 é o do \textbf{próprio estudante}.\\
\cd{.uni-channel-header\_\_cover} usa a primitiva
  \cd{--uni-color-neutral-50}, não um token semântico & No escuro vira uma laje
  branca.\\
\zebra \arq{styles-unootube.css} repete o \cd{@import} da Roboto no
  \textbf{meio} do arquivo & O Vite descarta com aviso; o do topo continua
  valendo, então a fonte carrega. É ruído, não bug.\\
Não existe componente de seleção (só \cd{TextField}, \cd{TextArea},
  \cd{PasswordInput}, \cd{SearchInput}) & \cd{<select>} do hub usa
  \cd{.unooria-select}, com os tokens do pacote --- mesma altura, raio e cores do
  \cd{TextField}.\\
\bottomrule
\end{tabularx}

\subsection{Lacunas abertas na 0.1.17 (exp-02 1.1.0)}

\begin{tabularx}{\linewidth}{@{}L{40mm}X@{}}
\toprule
\headrow \thd{Lacuna} & \thd{Consequência}\\
\midrule
\cd{unootube} não publica \cd{Tag} nem \cd{Badge} & Nenhum dos dois arquivos de
  origem desenha um. A etiqueta \enquote{LIMITE ATINGIDO} do diálogo de limite é
  compensação local --- e sai \textbf{neutra}, não âmbar como o Figma pede: a
  seção não tem token de aviso, e derivar um seria inventar valor.\\
\zebra O véu do \cd{ShortsPlayer} tem altura \textbf{fixa} de 194\,px & Ele
  existe para a legenda ficar legível. Numa miniatura de 88\,px sem legenda
  (\cd{showChannel} \cd{false}), o véu é mais alto que a miniatura inteira e só a
  faixa mais escura do gradiente aparece: o sintoma é \textbf{retângulo preto no
  lugar do vídeo}, com o vídeo carregado por baixo. Compensado com
  \cd{display: none} nas miniaturas da bandeja e do diálogo. Deveria ser
  proporcional, ou preso ao \cd{showChannel}.\\
\cd{Button} do clone não tem variante de ênfase & \enquote{Entendi!} e
  \enquote{Salvar} recebem preenchimento por ligação local, com os tokens da
  própria marca.\\
\bottomrule
\end{tabularx}

\subsection{Lacunas abertas em \cd{unooria} na 0.1.17 (hub, 2026-09-05)}

As duas lacunas abaixo apareceram nas telas de turma do professor, onde a
listagem tem uma linha por estudante --- 40 numa sala cheia. As duas são
compensadas em \arq{packages/ui/styles/plataforma.css} e viraram pedido no
\arq{storybook-unooria}.

\begin{tabularx}{\linewidth}{@{}L{40mm}X@{}}
\toprule
\headrow \thd{Lacuna} & \thd{Consequência}\\
\midrule
\cd{Table} não tem cabeçalho fixo & \cd{TableProps} não publica
  \cd{stickyHeader} nem \cd{maxHeight}, e o \cd{.unooria-table-container} só rola
  na horizontal. Numa turma de 40 estudantes o cabeçalho sai da tela junto com a
  página, e o professor lê \enquote{sim $\cdot$ --- $\cdot$ 3} sem saber de qual
  coluna. Compensado pela classe \cd{.tabela-rolavel}, que dá altura máxima ao
  contêiner (tela útil menos o cromo, piso de 20\,rem) e prende o \cd{<th>} com
  \cd{position: sticky}. Vale para as duas famílias de tabela do hub: a
  \cd{<table>} crua e a do DS (\cd{className="tabela-rolavel"}).\\
\zebra \cd{Button} não tem variante de perigo discreta & \cd{ButtonVariant} é
  \cd{primary | secondary | outline | ghost | danger}: as três primeiras saem no
  navy da marca, que não diz \enquote{isto apaga}, e a \cd{danger} é sólida ---
  peso de ação primária, errado para \enquote{Zerar a turma} ao lado da leitura
  da aula. Compensado por \cd{.acao-destrutiva}, que troca o navy do
  \cd{outline} pelo token de perigo e mantém o fundo transparente.\\
\bottomrule
\end{tabularx}

\begin{atencao}
\textbf{Variante que não existe não é erro de compilação --- é botão sem cor.} O
\cd{Button} monta a classe por interpolação (\cd{unooria-button--variant-\$\{variant\}}),
então um \cd{variant="tertiary"} produz uma classe que nenhuma folha define: o
botão perde fundo, cor e borda e fica com os estilos de base. Era o que
acontecia com \enquote{Zerar a turma} do UnooGPT e do Eco no Grupo, e com o
\enquote{confirmar} do diálogo (\cd{variant="destructive"}). O \emph{typecheck}
não pega. Ao escrever um \cd{variant}, confira contra \cd{ButtonVariant}.
\end{atencao}

\subsection{Lacunas abertas em \cd{whatsappmobile} na 0.1.17 (exp-07 0.1.0)}

A seção tem 29 componentes e cobre o cromo inteiro do mensageiro --- barra de
status, cabeçalho de conversa, balão com mídia e documento, campo de escrever,
diálogo, folha inferior. As quatro peças que a exp-07 pede e ela não publica
estão abaixo; \textbf{três saem de composição}, e só uma é compensação de
verdade (\adrr{ADR-68}).

\begin{tabularx}{\linewidth}{@{}L{40mm}X@{}}
\toprule
\headrow \thd{Lacuna} & \thd{Como está resolvido}\\
\midrule
\cd{ForwardedBadge} --- o selo \enquote{Encaminhada com frequência} & Composição:
  \cd{Icon name="forward"} (o pacote tem o glifo) com os tokens de texto
  secundário. Nenhuma cor inventada.\\
\zebra \cd{PinnedMessageBanner} --- a barra do alfinete & Composição sobre a
  superfície de \emph{appbar}. O pacote tem o \emph{tone} \cd{pin}, com a cor
  certa, mas \textbf{nenhum glifo de alfinete} --- ele é desenhado localmente,
  junto de \cd{play} e \cd{pause}, que também faltam.\\
\cd{ReplyQuoteSelectionSheet} --- o menu das 3 respostas & Composição com
  \cd{Radio} e um card de citação. \textbf{Não use \cd{ListItem}}: o título dele
  tem \cd{white-space: nowrap} --- certo numa lista de conversas, e errado aqui,
  porque truncar as opções faria a criança escolher sem ler.\\
\zebra \cd{ChatAudioBubbleWithTranscript} --- o balão de voz &
  \textbf{A única compensação real.} \cd{MessageBubble} aceita texto,
  \cd{.Media} (imagem) e \cd{.Document}; não há player de áudio na seção. O
  balão local traz play, onda em \emph{divs} (sem Canvas, por causa do
  \emph{framerate} em Chromebook escolar) e a transcrição sempre visível.\\
\bottomrule
\end{tabularx}

\begin{atencao}
\textbf{A seção publica só o tema escuro} (\cd{WaTheme = 'dark'}; o CSS traz
\cd{:root, [data-wa-theme=dark]} e nada mais). Os wireframes da exp-07 descrevem
telas claras; derivar a paleta é proibido, pela mesma regra do \cd{unitube}. O
contraste do escuro (\cd{\#D3DAE0} sobre \cd{\#101D25}, cerca de 11:1) passa o
AA com folga --- ver \adrr{ADR-68}.

E o \cd{PhoneFrame} do pacote \textbf{não entra} na experiência: quem desenha o
aparelho é o hub, pelo \cd{viewport} do manifesto. Empilhar os dois daria um
celular dentro de outro.
\end{atencao}

\subsection{Avisos que ficaram obsoletos}

\begin{nota}
Dois avisos herdados do repositório de origem estão \textbf{obsoletos na 0.1.11},
medidos no pacote:

\begin{itemize}
  \item o \cd{MetricCard} \textbf{não inventa mais tendência} --- a prop é
    \cd{tagText?} e a documentação diz \enquote{sem ela a etiqueta não é
    renderizada};
  \item o \cd{PositiveNegativeBarChart} tem escala \textbf{fixa} em $\pm 100$, que
    é exatamente o que os eixos do circumplexo pedem --- então a barra divergente
    \emph{não} é mais escrita à mão.
\end{itemize}

O terceiro segue valendo: \cd{AlertBox} é \textbf{diálogo} de 436\,px; para faixa
de contexto, o componente é \cd{FeedbackAlert} (e a prop é \cd{type}, não
\cd{variant}).
\end{nota}

\begin{atencao}
\textbf{Sobrescrever classe do pacote exige DOIS seletores}
(\cd{.app-shell > .x}): o CSS do pacote é injetado pelo módulo do componente,
ou seja, \emph{depois} do CSS da experiência no grafo de importação. Com peso
igual, quem ganha é sempre o pacote.
\end{atencao}

\section{Cor da plataforma: classe escrita à mão}

\begin{atencao}
A cor das superfícies da plataforma vem de classes escritas à mão em
\arq{packages/ui/styles/plataforma.css}, e o tema do Tailwind está
\textbf{vazio} (\cd{theme.extend: \{\}}) --- de propósito, para que nada da
marca institucional vaze para a superfície do estudante, nem por nome de classe.

A consequência é uma armadilha real: \textbf{classe de cor que não existe ali não
é erro em lugar nenhum}. Em texto, ela herda a cor e passa desapercebida; em
\textbf{fundo}, o elemento fica branco no branco. Foi o que aconteceu com
\cd{bg-unooria-neutral-900}: a letra da bolha ficou invisível. O \emph{smoke} da
exp-70 checa o \cd{backgroundColor} \textbf{computado} do círculo por isso ---
\cd{toBeVisible()} não pega.

Existem hoje apenas:
\cd{bg-unooria-\{navy-800, neutral-50, neutral-100, neutral-200\}},
\cd{border-unooria-neutral-200} e
\cd{text-unooria-\{danger, navy-900, neutral-500, neutral-700, neutral-800,
neutral-950\}}.
\end{atencao}

\begin{nota}
O \emph{preflight} do Tailwind fica \textbf{ligado}. Desligá-lo parecia inofensivo
--- o pacote traz o próprio CSS --- mas o pacote estiliza \textbf{componentes},
não elementos: sem o \emph{reset}, todo \cd{<a>} volta a ser sublinhado e todo
\cd{<ul>} volta a ter marcador. Foi o que o \emph{smoke} visual pegou nos
\emph{cards} de produto.
\end{nota}

\section{Exceção registrada}

\begin{decisao}
\textbf{A exp-70 está fora da regra do \emph{design system}}, deliberadamente:
é jogo em \emph{canvas} 3D, sem DOM de interface (\adrr{ADR-18}). A pendência da
exp-02 --- que tinha CSS próprio em \arq{src/estilos} --- foi
\textbf{resolvida}: ela consome a seção \cd{unitube} pela ponte
\cd{@edu/ui/ds-unootube}, com a aparência escolhida pelo estudante.
\end{decisao}
